\documentclass[11pt,a4paper]{article}

\usepackage[T1]{fontenc}
\usepackage{lmodern}
\usepackage{geometry}
\usepackage{amsmath,amssymb}
\usepackage{booktabs}
\usepackage{tabularx}
\usepackage{array}
\usepackage{float}
\usepackage{graphicx}
\usepackage{hyperref}
\usepackage{xcolor}
\usepackage{enumitem}
\usepackage{caption}

\geometry{left=2.5cm,right=2.5cm,top=2.5cm,bottom=2.5cm}
\hypersetup{
  colorlinks=true,
  linkcolor=blue!55!black,
  citecolor=blue!55!black,
  urlcolor=blue!55!black
}
\setlist[itemize]{leftmargin=2em}
\setlist[enumerate]{leftmargin=2em}
\captionsetup{font=small,labelfont=bf}
\emergencystretch=3em

\newcommand{\ccac}{\textsc{CCAC}}
\newcommand{\ours}{\textsc{Focal+Soft}}
\newcommand{\ind}{\mathrm{IND}}
\newcommand{\ood}{\mathrm{OOD}}
\newcommand{\qc}{Q_c}

\title{\textbf{A Small-Scale Extension of CCAC with Improved OOD Classification}\\
\large RL Course Final Report}
\author{
\begin{tabular}{ll}
Name: Wang Xianghong & Student ID: 250010179 \\
Name: Lin Xinying & Student ID: 250010029 \\
Name: Hu Ronghuai & Student ID: 250010030 \\
Name: Wang Zihao & Student ID: 525180
\end{tabular}
}
\date{June 2026}

\begin{document}

\maketitle

\begin{abstract}
Offline safe reinforcement learning aims to learn high-reward policies from a
fixed dataset while satisfying safety constraints. Constraint-Conditioned
Actor-Critic (\ccac) addresses this problem by conditioning the actor and
critics on the remaining safety budget, allowing one trained policy to be
evaluated under multiple target costs. This report presents a course-project
scale extension of an existing \ccac{} implementation. The extension replaces
the original binary cross-entropy loss for the OOD classifier with focal loss
and uses the classifier output as a soft weight in the OOD cost-critic update.
The evaluation follows a three-level pipeline. Test A trains only the OOD
classifier and measures unsafe/OOD false negatives. Test B trains a simplified
cost critic and measures matched-budget OOD-vs-IND cost separation. Test C runs
50k-update full policy training and evaluates reward and realized cost under
multiple target costs. The results show stable mechanism-level improvements:
the proposed variant substantially reduces false negative rates on BallRun and
CarRun and improves matched-budget cost-critic separation on CarRun across two
seeds. However, full policy results are mixed. The variant improves reward or
strict-budget safety in several BallRun settings, but it fails on BallRun seed
2 at target cost 1 and does not solve CarRun policy safety. Therefore, the
\ours{} variant should be viewed as a useful diagnostic and small-scale
extension rather than a robust new offline safe RL algorithm.
\end{abstract}

\noindent\textbf{Keywords:} offline safe reinforcement learning; CCAC; OOD
detection; focal loss; cost critic; constraint-conditioned policy

\section*{Member Contributions}

The project work was divided into four main parts, with one member taking the
main responsibility for each part.

\begin{table}[H]
  \centering
  \caption*{Main responsibilities}
  \begin{tabularx}{0.9\textwidth}{clX}
    \toprule
    No. & Project part & Responsible member \\
    \midrule
    1 & Project material collection & Lin Xinying \\
    2 & Code implementation & Wang Zihao \\
    3 & Presentation delivery & Wang Xianghong \\
    4 & Experimental report writing & Hu Ronghuai \\
    \bottomrule
  \end{tabularx}
\end{table}

The entries above describe the main stage each student was responsible for.
Beyond these primary responsibilities, all group members participated in idea
discussion, project iteration, result interpretation, and final refinement. The
overall contribution ratio is therefore considered equal across the four
members, namely 25\% per member.

\section{Introduction}

Safe reinforcement learning optimizes task reward while controlling cumulative
safety cost. In the offline setting, the learner is restricted to a fixed
dataset and cannot interact with the environment to correct high-risk actions.
This makes distribution shift especially problematic: when the learned policy
selects actions outside the support of the offline data, value estimation can
become unreliable and safety constraints may be violated during evaluation.

\ccac{} is a recent offline safe RL method that learns a
constraint-conditioned actor-critic policy \cite{guo2025ccac}. It treats the
remaining safety budget \(\kappa_t\) as part of the input to the policy,
critics, and generative model. As a result, the same trained policy can be
evaluated zero-shot under different target cost thresholds. In the original
\ccac{} design, a conditional variational autoencoder (CVAE) generates
budget-conditioned state-action samples, an OOD classifier identifies generated
out-of-distribution samples, and the cost critic receives conservative updates
on risky generated samples.

This project focuses on a small but safety-relevant issue in the \ccac{}
pipeline. If the OOD classifier misses unsafe or out-of-distribution samples,
the cost critic may not receive enough conservative safety signal, and the
final policy may remain unsafe. The original implementation trains the
classifier with binary cross entropy (BCE) and uses a hard threshold to select
OOD samples. A hard OOD decision treats near-threshold samples very
differently, for example assigning different behavior to scores 0.49 and
0.51. In safety-critical settings, false negatives are particularly dangerous:
an unsafe OOD sample classified as in-distribution may avoid conservative cost
regularization.

Motivated by this observation, this report compares original \ccac{} with one
selected variant:
\begin{itemize}
  \item Baseline: original \ccac{}, with \texttt{classifier\_loss=bce} and
  \texttt{ood\_mode=hard}.
  \item Ours: \ours{}, with \texttt{classifier\_loss=focal},
  \texttt{ood\_mode=soft}, \texttt{focal\_alpha=0.75}, and
  \texttt{focal\_gamma=2.0}.
\end{itemize}

The contribution of this project is threefold. First, it implements two
lightweight modifications to \ccac{}: a focal-loss OOD classifier and a soft
OOD weighting rule for the cost-critic update. Second, it evaluates the change
through a three-level validation pipeline that separates classifier behavior,
cost-critic diagnostics, and final policy performance. Third, it reports both
positive and negative evidence. The mechanism-level results are consistently
better, but the policy-level results are not uniformly safer.

\section{Background}

\subsection{Offline Safe Reinforcement Learning}

Offline safe RL can be formulated as a constrained optimization problem:
\begin{equation}
  \max_{\pi} \ \mathbb{E}_{\pi}\left[R(\tau)\right],
  \quad
  \mathrm{s.t.}\quad
  \mathbb{E}_{\pi}\left[C(\tau)\right] \le \epsilon,
  \label{eq:offline-safe-rl}
\end{equation}
where \(R(\tau)\) is the trajectory return, \(C(\tau)\) is the cumulative
safety cost, and \(\epsilon\) is the target cost threshold. Since the training
dataset is fixed, the learned policy cannot safely explore the environment.
This makes distributional shift and value extrapolation central challenges.
The DSRL benchmark \cite{liu2024dsrl} provides a suite of offline safe RL tasks
and reports reward, cost, and normalized scores for systematic evaluation.

\subsection{Review of CCAC}

\ccac{} conditions policy learning on the remaining safety budget. Given
budget \(\kappa_t\), the policy is written as
\begin{equation}
  \pi_\theta(a_t \mid s_t, \kappa_t),
  \qquad
  \kappa_{t+1} = \kappa_t - c_t.
  \label{eq:budget-update}
\end{equation}
This design allows one actor to adapt to multiple target cost thresholds at
evaluation time.

The algorithm consists of four main components. A CVAE generates
budget-conditioned state-action samples. An OOD classifier identifies generated
samples that are outside the data distribution. Reward and cost critics
estimate return and safety cost. The actor is trained with a Lagrangian-style
objective that balances reward maximization and predicted cost control. For
safety, the interaction between the OOD classifier and the cost critic is
especially important, because conservative cost updates on generated samples
shape how the learned policy avoids risky actions outside the dataset support.

\section{Method}

\subsection{Focal OOD Classifier}

The original implementation trains the OOD classifier using BCE. This project
replaces it with focal loss \cite{lin2017focal}:
\begin{equation}
  \mathrm{FL}(p_t)
  =
  \alpha_t (1-p_t)^\gamma \mathrm{BCE}(p,y),
  \label{eq:focal}
\end{equation}
where \(p\) is the classifier's predicted OOD probability, \(y\in\{0,1\}\) is
the binary label, \(p_t=p\) when \(y=1\), and \(p_t=1-p\) otherwise. The
experiment uses \(\alpha=0.75\) and \(\gamma=2.0\). Focal loss down-weights
easy examples and gives relatively more weight to hard examples. In this
project, the goal is to reduce false negatives, namely unsafe/OOD samples that
are incorrectly treated as in-distribution.

\subsection{Soft OOD Weighting}

The original hard OOD mode uses a threshold:
\begin{equation}
  w_{\ood}(s,a,\kappa)
  =
  \mathbb{I}\left[h_\psi(s,a,\kappa) \ge \tau\right],
  \label{eq:hard-ood}
\end{equation}
where \(h_\psi\) is the OOD classifier and \(\tau=0.5\). The proposed soft mode
uses the classifier output directly as a weight:
\begin{equation}
  w_{\ood}(s,a,\kappa)
  =
  \mathrm{clamp}\left(h_\psi(s,a,\kappa),0,1\right).
  \label{eq:soft-ood}
\end{equation}
This avoids a discontinuous decision around a single threshold. Suspicious
samples that do not exceed the hard threshold can still contribute to the OOD
cost signal.

\subsection{Effect on Cost-Critic Learning}

In \ccac{}, generated OOD samples are used to impose a conservative
Lagrangian-style constraint on the cost critic. With soft OOD weighting, the
OOD cost loss can be interpreted as
\begin{equation}
  \mathcal{L}_{\ood}
  =
  - w_{\ood}(s,a,\kappa)\lambda_\phi(\kappa)
  \left(\qc(s,a,\kappa)-\kappa\right).
  \label{eq:ood-cost-loss}
\end{equation}
Samples with higher OOD scores receive stronger conservative updates. Samples
near the threshold are no longer discarded entirely.

\section{Experimental Setup}

\subsection{Tasks and Compared Methods}

The experiments use DSRL offline safe RL tasks. The main task is
\texttt{OfflineBallRun-v0}, which is used for full policy comparison over
seeds \(0/1/2\). The second task is \texttt{OfflineCarRun-v0}, which is used
for repeated mechanism-level experiments over seeds \(0/1\) and one seed-0
full-policy check. The CarRun dataset is a locally supplied
\texttt{SafetyCarRun-v0-40-651.hdf5} file with 130200 transitions. A planned
BallCircle fallback task was not included because the public dataset server
returned HTTP 502 during the final expansion.

Table~\ref{tab:protocol} summarizes the three-level validation pipeline. Test
A and Test B do not replace full policy evaluation. Instead, they isolate
whether the classifier and cost critic change in the intended direction before
the method is evaluated as a full policy.

\begin{table}[H]
  \centering
  \caption{Experiment protocol and training budget}
  \label{tab:protocol}
  \begin{tabularx}{\textwidth}{llclX}
    \toprule
    Test & Task & Seeds & Steps & Purpose \\
    \midrule
    Test A & BallRun, CarRun & BallRun 0; CarRun 0/1 & 5000 &
    OOD classifier diagnostics, focusing on FNR. \\
    Test B & BallRun, CarRun & BallRun 0; CarRun 0/1 & 10000 &
    Matched-budget cost-critic separation. \\
    Test C & BallRun & 0/1/2 & 50000 &
    Full policy comparison. \\
    Test C & CarRun & 0 & 50000 &
    Second-task policy-level check. \\
    \bottomrule
  \end{tabularx}
\end{table}

\subsection{Metrics}

Test A reports AUROC, AUPRC, false negative rate (FNR), false positive rate
(FPR), and expected calibration error (ECE). The most important safety-side
metric is FNR, because it measures how often unsafe/OOD samples are incorrectly
treated as in-distribution.

Test B reports three cost-critic diagnostics:
\begin{itemize}
  \item \texttt{qc\_matched\_gap}: mean OOD \(\qc\) minus mean matched IND
  \(\qc\);
  \item \texttt{qc\_selected\_matched\_gap}: mean selected OOD \(\qc\) minus
  mean matched IND \(\qc\);
  \item \texttt{ood\_selected\_rate}: the fraction of generated samples selected
  or effectively weighted as OOD.
\end{itemize}

Test C uses \texttt{eval\_ccac.py} to evaluate target costs
\([1,2,5,10]\), with 20 episodes per target cost. A target is counted as safe
if \(\mathrm{real\_cost}\le\mathrm{target\_cost}\). The ``safe targets''
column reports how many of the four target costs satisfy this criterion.

\subsection{Environment}

The experiment records use Python 3.8, PyTorch 2.4.1+cu121, Gymnasium, DSRL,
and OSRL. The main GPU is an NVIDIA H100 with 80GB HBM3 memory. Training uses
\texttt{WANDB\_MODE=offline}, and the local HDF5 dataset directory is specified
through \texttt{DSRL\_DATASET\_DIR}.

\section{Results}

\subsection{Test A: OOD Classifier}

Table~\ref{tab:test-a} shows the classifier-only diagnostics. The \ours{}
variant substantially reduces FNR in all three settings. On BallRun seed 0,
FNR decreases from 0.10380 to 0.03565. On CarRun seed 0, it decreases from
0.13206 to 0.05261. On CarRun seed 1, it decreases from 0.18108 to 0.06185.
This supports the first mechanism hypothesis: focal loss makes the classifier
miss fewer unsafe/OOD samples.

\begin{table}[H]
  \centering
  \caption{Test A: OOD classifier metrics}
  \label{tab:test-a}
  \resizebox{\textwidth}{!}{
  \begin{tabular}{llrrrrrr}
    \toprule
    Task & Method & Seed & AUROC & AUPRC & FNR & FPR & ECE \\
    \midrule
    BallRun & original & 0 & 0.94270 & 0.93778 & 0.10380 & 0.17205 & 0.02227 \\
    BallRun & \ours & 0 & 0.93943 & 0.93326 & 0.03565 & 0.28166 & 0.10002 \\
    CarRun & original & 0 & 0.92558 & 0.88536 & 0.13206 & 0.17275 & 0.03152 \\
    CarRun & \ours & 0 & 0.92649 & 0.88579 & 0.05261 & 0.28111 & 0.13474 \\
    CarRun & original & 1 & 0.92868 & 0.89186 & 0.18108 & 0.12899 & 0.01671 \\
    CarRun & \ours & 1 & 0.92992 & 0.89303 & 0.06185 & 0.25883 & 0.12755 \\
    \bottomrule
  \end{tabular}}
\end{table}

The improvement is not free. Compared with original \ccac{}, \ours{} has a
higher FPR and worse ECE. This means that the classifier becomes more
conservative but less calibrated. Therefore, the correct interpretation is a
safety-oriented trade-off rather than a uniformly better classifier.

\subsection{Test B: Matched-Budget Cost-Critic Separation}

Test B asks whether generated OOD samples receive higher cost estimates under
the same query budget. The results are shown in Table~\ref{tab:test-b}. On
CarRun seed 0 and seed 1, \ours{} improves both
\texttt{qc\_matched\_gap} and \texttt{qc\_selected\_matched\_gap}, suggesting
that the more conservative classifier and soft OOD weighting strengthen the OOD
cost signal. On BallRun, original \ccac{} has a negative
\texttt{qc\_matched\_gap}, while \ours{} produces a positive matched gap. This
indicates that the hard-mask baseline can have unstable cost-critic scale in
this small experiment.

\begin{table}[H]
  \centering
  \caption{Test B: matched-budget cost-critic separation}
  \label{tab:test-b}
  \resizebox{\textwidth}{!}{
  \begin{tabular}{llrrrr}
    \toprule
    Task & Method & Seed & \texttt{qc\_matched\_gap} &
    \texttt{qc\_selected\_matched\_gap} & \texttt{ood\_selected\_rate} \\
    \midrule
    BallRun & original & 0 & -26.829 & 88.324 & 0.4186 \\
    BallRun & \ours & 0 & 8.711 & 6.131 & 0.7555 \\
    CarRun & original & 0 & 2.641 & 4.929 & 0.3106 \\
    CarRun & \ours & 0 & 3.632 & 6.644 & 0.4868 \\
    CarRun & original & 1 & 3.267 & 6.069 & 0.2267 \\
    CarRun & \ours & 1 & 5.031 & 8.027 & 0.4255 \\
    \bottomrule
  \end{tabular}}
\end{table}

\ours{} also has a higher \texttt{ood\_selected\_rate}, which is consistent
with its more conservative classifier behavior in Test A. Overall, Test B
supports the second mechanism hypothesis: soft OOD weighting improves
matched-budget OOD-vs-IND cost separation.

\subsection{Test C: Full Policy Evaluation}

The full policy results are more complicated. Table~\ref{tab:test-c}
summarizes average reward, average realized cost, and safe targets across the
four evaluation budgets.

\begin{table}[H]
  \centering
  \caption{Test C: full policy summary}
  \label{tab:test-c}
  \begin{tabular}{llrrrl}
    \toprule
    Task & Method & Seed & Avg reward & Avg real cost & Safe targets \\
    \midrule
    BallRun & original & 0 & 347.274 & 0.000 & 4/4 \\
    BallRun & \ours & 0 & 422.771 & 0.000 & 4/4 \\
    BallRun & original & 1 & 464.176 & 19.000 & 3/4 \\
    BallRun & \ours & 1 & 420.406 & 3.250 & 4/4 \\
    BallRun & original & 2 & 409.096 & 0.000 & 4/4 \\
    BallRun & \ours & 2 & 443.294 & 9.650 & 3/4 \\
    CarRun & original & 0 & 859.456 & 181.200 & 0/4 \\
    CarRun & \ours & 0 & 913.301 & 183.213 & 0/4 \\
    \bottomrule
  \end{tabular}
\end{table}

On BallRun, \ours{} provides useful but not uniform improvements. In seed 0,
it improves reward under all target costs while keeping zero realized cost. In
seed 1, it fixes a strict-budget failure: original \ccac{} has real cost 74 at
target cost 1, while \ours{} has real cost 0. However, it loses reward at
target costs 5 and 10. In seed 2, \ours{} improves reward at target costs 2 and
5 while remaining safe there, but it fails badly at target cost 1 with real
cost 33.6.

CarRun is the clearest negative policy result. Table~\ref{tab:carrun-policy}
shows the key numbers. \ours{} obtains higher reward than original \ccac{}, but
realized cost does not improve. Both methods exceed all target cost thresholds
by a large margin.

\begin{table}[H]
  \centering
  \caption{Key CarRun policy results}
  \label{tab:carrun-policy}
  \begin{tabular}{lrrrr}
    \toprule
    Method & Target 1 cost & Target 5 cost & Avg reward & Safe targets \\
    \midrule
    original & 181.00 & 181.25 & 859.456 & 0/4 \\
    \ours & 183.05 & 183.55 & 913.301 & 0/4 \\
    \bottomrule
  \end{tabular}
\end{table}

This result is important because it shows that better mechanism-level OOD
diagnostics do not automatically guarantee safer policy behavior.

\section{Discussion}

The main finding is the gap between mechanism-level improvements and
policy-level behavior. Test A and Test B show that \ours{} changes the internal
OOD safety signal in the intended direction. The classifier misses fewer
unsafe/OOD samples, and the cost critic assigns higher cost estimates to OOD
samples under matched budgets. This supports the design rationale behind focal
loss and soft OOD weighting.

However, the policy-level results are not stable. In some BallRun settings,
\ours{} increases reward or fixes strict-budget safety failures. In other
settings, especially BallRun seed 2 at target cost 1, it introduces a clear
safety failure. On CarRun, both original \ccac{} and \ours{} remain unsafe
after 50k updates. This suggests that the lightweight OOD modification is not
sufficient to solve budget-conditioned safe control on harder tasks.

The three-level evaluation pipeline is therefore useful. If only Test A and
Test B were reported, the result would look more positive than it should. If
only Test C were reported, it would be difficult to determine whether failures
came from the classifier, cost critic, or actor training. Separating mechanism
diagnostics from full policy evaluation makes the final conclusion more
interpretable.

\section{Limitations}

This project has several limitations:
\begin{itemize}
  \item Full policy training uses only 50k updates, which is much smaller than a
  full benchmark-scale reproduction.
  \item Policy-level evaluation covers only BallRun seeds \(0/1/2\) and CarRun
  seed 0.
  \item Focal loss reduces FNR but increases FPR and worsens ECE, meaning that
  classifier calibration becomes worse.
  \item Soft OOD weighting improves the mechanism-level OOD cost signal, but it
  does not solve the CarRun policy safety failure.
  \item The project does not systematically compare against other offline safe
  RL or offline RL baselines such as CQL \cite{kumar2020cql} and CDT
  \cite{liu2023cdt}.
  \item The BallCircle fallback task was excluded because the public dataset
  server failed during the final expansion.
\end{itemize}

\section{Reproducibility Notes}

The main experiment scripts are under
\path{CCAC/validation_experiments/scripts/}. Before running experiments, the
environment variables \texttt{PYTHONPATH}, \texttt{WANDB\_MODE=offline}, and
\texttt{DSRL\_DATASET\_DIR} should be set. The main entry points are:
\begin{itemize}
  \item \path{run_test_a_classifier.sh}: classifier-only diagnostics;
  \item \path{run_test_b_cost_critic.sh}: cost-critic separation;
  \item \path{train_original_ccac_50k.sh}: original \ccac{} 50k baseline;
  \item \path{OSRL/examples/train/train_ccac.py}: full policy training;
  \item \path{OSRL/examples/eval/eval_ccac.py}: multi-target-cost policy
  evaluation.
\end{itemize}
Full commands, run directories, and raw outputs are recorded in
\path{CCAC/validation_experiments/commands.md} and
\path{CCAC/validation_experiments/results.md}. Report-ready tables are recorded
in \path{CCAC/validation_experiments/final_report_results.md}.

\section{Conclusion}

This report presents a small-scale extension of \ccac{} that trains the OOD
classifier with focal loss and uses soft OOD scores in the cost-critic update.
The proposed variant improves mechanism-level diagnostics: it reduces unsafe
OOD false negatives and improves matched-budget cost-critic separation.
However, full policy behavior remains mixed. The method improves several
BallRun settings but also has a clear strict-budget failure, and it does not
solve CarRun safety.

The appropriate conclusion is therefore cautious. \ours{} is a useful
diagnostic and small-scale extension, not a robust new offline safe RL
algorithm. Future work should evaluate the idea on more DSRL tasks, more seeds,
longer training budgets, and stronger baselines, while also studying the
relationship between classifier calibration, OOD cost regularization strength,
and policy stability.

\bibliographystyle{ieeetr}
\bibliography{references}

\end{document}
